% ============================================================
%  Beamer Presentation Template
%  Topic: RL-based Scheduling for Data Centers
% ============================================================
\documentclass[aspectratio=169]{beamer}  % 16:9 widescreen (like PowerPoint default)

% ---------- Theme -------------------------------------------
\usetheme{Madrid}
\usecolortheme{default}
% Other popular themes: AnnArbor, Berlin, Boadilla, CambridgeUS,
%   Copenhagen, Darmstadt, Dresden, Frankfurt, Hannover, Luebeck,
%   Malmoe, Montpellier, PaloAlto, Pittsburgh, Rochester, Singapore,
%   Szeged, Warsaw
% Other colour themes: albatross, beaver, beetle, crane, dolphin,
%   dove, fly, lily, orchid, rose, seagull, seahorse, spruce, whale, wolverine

% ---------- Packages ----------------------------------------
\usepackage{graphicx}        % Images
\usepackage{booktabs}        % Nice tables (\toprule, \midrule, \bottomrule)
\usepackage{amsmath}         % Math
\usepackage{hyperref}        % Clickable links
\usepackage{tikz}            % Diagrams (optional)
\usepackage{algorithm2e}     % Pseudocode (optional)

% ---------- Title info --------------------------------------
\title[RL Scheduling for DCs]{RL-based Scheduling for Data Centers}
\subtitle{A Reinforcement Learning Approach}
\author[A. Diaz Gonzalez]{Abel Diaz Gonzalez}
\institute[Your University]{%
  Department / Lab Name \\
  Your University or Institution
}
\date{June 2026}
\logo{\includegraphics[height=0.6cm]{logo}}  % remove or replace with your logo file

% ============================================================
\begin{document}

% ---- Title slide -------------------------------------------
\begin{frame}
  \titlepage
\end{frame}

% ---- Table of contents -------------------------------------
\begin{frame}{Outline}
  \tableofcontents
\end{frame}

% ============================================================
\section{Introduction}
% ============================================================

\begin{frame}{Motivation}
  \begin{itemize}
    \item Data centers consume a significant portion of global energy.
    \item Efficient job scheduling reduces cost and improves throughput.
    \item Traditional heuristics (FIFO, SJF) are sub-optimal under dynamic workloads.
    \item<2-> \textbf{Reinforcement Learning} can learn adaptive scheduling policies.
  \end{itemize}

  \vspace{0.5em}
  \begin{block}{Research Question}
    Can an RL agent learn a scheduling policy that outperforms hand-crafted
    heuristics in heterogeneous data center environments?
  \end{block}
\end{frame}

% ---- Two-column slide example ------------------------------
\begin{frame}{Problem Statement}
  \begin{columns}[T]
    \column{0.5\textwidth}
      \textbf{Challenges}
      \begin{itemize}
        \item Large, dynamic action space
        \item Partial observability
        \item Multi-objective optimisation
        \item Resource heterogeneity
      \end{itemize}

    \column{0.5\textwidth}
      \textbf{Goals}
      \begin{itemize}
        \item Minimise job completion time
        \item Maximise resource utilisation
        \item Reduce energy consumption
        \item Handle bursty workloads
      \end{itemize}
  \end{columns}
\end{frame}

% ============================================================
\section{Background}
% ============================================================

\begin{frame}{Reinforcement Learning Basics}
  \begin{columns}
    \column{0.55\textwidth}
      \begin{itemize}
        \item \textbf{Agent} interacts with an \textbf{Environment}.
        \item Observes \textbf{state} $s_t$, takes \textbf{action} $a_t$.
        \item Receives \textbf{reward} $r_t$ and transitions to $s_{t+1}$.
        \item Goal: maximise cumulative reward
          \[
            G_t = \sum_{k=0}^{\infty} \gamma^k r_{t+k}
          \]
      \end{itemize}

    \column{0.4\textwidth}
      % Replace with a real figure if available
      \begin{figure}
        \centering
        \fbox{\rule{0pt}{2.5cm}\rule{4cm}{0pt}}  % placeholder box
        \caption{Agent--Environment loop}
      \end{figure}
  \end{columns}
\end{frame}

% ---- Table slide example -----------------------------------
\begin{frame}{Related Work Comparison}
  \begin{table}
    \centering
    \begin{tabular}{lccc}
      \toprule
      \textbf{Method} & \textbf{Adaptive} & \textbf{Multi-resource} & \textbf{Scalable} \\
      \midrule
      FIFO            & \xmark & \xmark & \cmark \\
      SJF             & \xmark & \cmark & \cmark \\
      Tetris          & \cmark & \cmark & \xmark \\
      DeepRM          & \cmark & \cmark & \cmark \\
      \textbf{Ours}   & \cmark & \cmark & \cmark \\
      \bottomrule
    \end{tabular}
    \caption{Comparison of scheduling approaches}
  \end{table}
\end{frame}

% ============================================================
\section{Methodology}
% ============================================================

\begin{frame}{System Architecture}
  \begin{figure}
    \centering
    \fbox{\rule{0pt}{4cm}\rule{10cm}{0pt}}  % replace with \includegraphics{arch}
    \caption{Overview of the proposed RL scheduling framework}
  \end{figure}
\end{frame}

\begin{frame}{MDP Formulation}
  \begin{description}
    \item[State $s_t$] Queue of pending jobs + resource utilisation matrix
    \item[Action $a_t$] Job–machine assignment at each scheduling step
    \item[Reward $r_t$] $-\text{(average job slowdown)}$ per time step
    \item[Policy $\pi_\theta$] Parameterised by a deep neural network
  \end{description}

  \vspace{0.5em}
  \begin{alertblock}{Key Design Choice}
    We use \textbf{Proximal Policy Optimisation (PPO)} to ensure
    stable policy updates in the non-stationary scheduling environment.
  \end{alertblock}
\end{frame}

% ============================================================
\section{Experiments}
% ============================================================

\begin{frame}{Experimental Setup}
  \begin{itemize}
    \item \textbf{Simulator:} custom discrete-event data center simulator
    \item \textbf{Workloads:} synthetic Poisson arrivals + Google cluster traces
    \item \textbf{Baselines:} FIFO, SJF, Tetris, DeepRM
    \item \textbf{Metrics:} average job completion time, resource utilisation, energy
    \item \textbf{Hardware:} 4 $\times$ NVIDIA A100, 512 GB RAM
  \end{itemize}
\end{frame}

\begin{frame}{Results}
  \begin{columns}
    \column{0.5\textwidth}
      \begin{figure}
        \centering
        \fbox{\rule{0pt}{3.5cm}\rule{5cm}{0pt}}
        \caption{Average job completion time}
      \end{figure}

    \column{0.5\textwidth}
      \begin{figure}
        \centering
        \fbox{\rule{0pt}{3.5cm}\rule{5cm}{0pt}}
        \caption{Resource utilisation}
      \end{figure}
  \end{columns}
\end{frame}

% ============================================================
\section{Conclusion}
% ============================================================

\begin{frame}{Conclusion \& Future Work}
  \begin{block}{Summary}
    \begin{itemize}
      \item Proposed an RL-based scheduler for heterogeneous data centers.
      \item Achieves X\% reduction in average job completion time over best baseline.
      \item Scales to clusters of $N$ machines and $M$ job types.
    \end{itemize}
  \end{block}

  \begin{block}{Future Work}
    \begin{itemize}
      \item Multi-agent scheduling (one agent per cluster rack)
      \item Integration with real Kubernetes/YARN clusters
      \item Energy-aware reward shaping
    \end{itemize}
  \end{block}
\end{frame}

% ---- Thank you / Q&A slide ---------------------------------
\begin{frame}
  \centering
  \Huge\textbf{Thank You!} \\[1em]
  \large Questions? \\[2em]
  \normalsize
  \href{mailto:your@email.com}{your@email.com} \\
  \url{https://github.com/yourhandle}
\end{frame}

% ---- References --------------------------------------------
\begin{frame}[allowframebreaks]{References}
  % Option A: manual list
  \begin{thebibliography}{9}
    \bibitem{mao2016resource}
      H. Mao et al.,
      \textit{Resource Management with Deep Reinforcement Learning},
      HotNets, 2016.
    % Add more entries here
  \end{thebibliography}

  % Option B: BibTeX — uncomment and point to your .bib file
  % \bibliographystyle{plain}
  % \bibliography{references}
\end{frame}

\end{document}
